\chapter{Exemples et Contre-exemples}

\begin{exercise}
  On observe que pour quelques entiers naturels $n$, on a $n! < 2^n$.
  Peut-on dire que cette propriété est vraie pour tout $n$ ?

  Qu'en est-il de $n! < 3^n$ ?

  Est-il vrai que $n! < 5^n$ pour tout entier naturel $n$ ?
\end{exercise}

\begin{exercise}
  La propriété
  \[
   \dfrac{1}{n!} > \dfrac{8^n}{(2n)!}
  \]
  est-elle vraie pour tout entier naturel $n$ ?
\end{exercise}

% \begin{exercise}
%     On observe pour quelques entiers positifs $n$ que l’expression $2^n + 3^n$ est divisible par $5$.
%
%     \textbf{Question :} Peut-on en conclure qu’elle est divisible par $5$ pour \emph{tous} les entiers $n>0$ ?
%     Justifiez par une recherche de contre-exemple ou, au contraire, par une preuve arithmétique.
% \end{exercise}

\begin{exercise}
 L'affirmation : "les diagonales d'un polygone sont toujours à l'intérieur du polygone" est-elle vraie ?
\end{exercise}

\begin{exercise}
 Pour quelles valeurs de $n$ a-t-on :
 \begin{enumerate}
  \item \[ \dfrac{1}{n} \leq 0,01 \]
  \item \[ 0,49 \leq \dfrac{n + 13}{2n + 25} \leq 0,51 \]
  \item \[ 2,5 \leq \frac{1}{0!} + \frac{1}{1!} + \frac{1}{2!} + \frac{1}{3!} + \cdots + \frac{1}{n!}  \leq 4 \]
 \end{enumerate}
\end{exercise}

\begin{exercise}
    Considérons le polynôme $P(n) = n^2 - n + 5$.
    \begin{itemize}
      \item Vérifiez que pour les premières valeurs de $n$, $P(n)$ est premier.
      \item \textbf{Discutez :} un \emph{grand} nombre d’exemples favorables suffirait-il à prouver que $P(n)$ est toujours premier ?
      \item Cherchez un entier $n$ rendant $P(n)$ composé.
    \end{itemize}
    Reprendre ces questions avec $P(n) = n^2 - n + 11$, $P(n) = 2^{2^n}+1$.

    Est-il vrai que si $p$ est un nombre premier, alors $2^p - 1$ est premier.
\end{exercise}

\begin{exercise}
    On étudie la formule
    \[
      1^3 + 2^3 + 3^3 + \dots + n^3 \;\stackrel{?}{=}\; \biggl(\frac{n(n+1)}{2}\biggr)^2.
    \]
    \begin{itemize}
      \item Calculez la somme pour $n=1,2,3,4,5,6$ et comparez avec $\left(\dfrac{n(n+1)}{2}\right)^2$.
      \item Ces quelques vérifications constituent-elles une démonstration ?
      \item Proposez, si possible, un \textbf{raisonnement} général (par récurrence ou par décomposition algébrique) pour prouver (ou infirmer) la formule.
    \end{itemize}
\end{exercise}

\begin{exercise}
    Considérons la suite définie par $a_n = n! + 1$.
    \begin{itemize}
      \item Calculez $a_1, a_2, a_3, a_4, a_5$. Certains de ces termes sont premiers.
      \item Peut-on affirmer que \emph{tous} les termes de cette forme sont premiers ?
      \item Soit vous exhibez un terme qui n’est pas premier, soit vous tentez de prouver qu’ils le sont tous.
    \end{itemize}
\end{exercise}

\begin{exercise}
    Dans un pentagone convexe, on peut mesurer la somme de ses angles intérieurs. En traçant plusieurs pentagones à main levée, on trouve souvent approximativement $540^\circ$.
    \begin{itemize}
      \item Expliquez pourquoi on ne peut pas se contenter d’exemples tracés et mesurés pour conclure.
      \item Donnez un argument géométrique \emph{global} prouvant que la valeur de la somme des angles d’un pentagone donne toujours $540^\circ$ (indice : triangulez la figure).
    \end{itemize}
\end{exercise}

\begin{exercise}
Dans un triangle $ABC$ quelconque, on peut essayer de démontrer expérimentalement que les médianes se coupent en un même point $G$.
\begin{itemize}
  \item Effectuez trois tracés différents d’un triangle et de ses médianes, constatez leur intersection.
  \item Montrez pourquoi la mesure ou l’observation graphique ne constitue pas une preuve.
  \item Proposez un argument géométrique global pour expliquer ce phénomène.

\end{itemize}
\end{exercise}

\begin{exercise}
    Vrai ou faux, justifiez : \\

    \begin{itemize}
      \item " Tout entier positif pair plus grand qu’un certain seuil s’écrit comme somme de deux nombres dont l’un est multiple de 3."
      \item " $3^{2n} + 1$ est toujours divisible par 10."
      \item " Le nombre de diagonales d'un polygone à $n$ côtés est  $\displaystyle\frac{n(n-3)}{2}$.". \\
      \emph{Indications : Considérer d'abord le cas des polygones réguliers, convexes puis le cas général.}
      \item " La somme de deux multiples de 4 est toujours un multiple de 8."
      \item " Tout quadrilatère convexe peut se paver entièrement avec un certain nombre de triangles identiques."
          \end{itemize}
\end{exercise}

\begin{exercise}
Proposez vous-même un énoncé \emph{nouveau} qui \textbf{semble} toujours vrai au vu de quelques cas, mais dont la validité générale n’est pas évidente.
\begin{itemize}
  \item Donnez quelques exemples numériquement ou géométriquement.
  \item Invitez le lecteur à \emph{soit} trouver un contre-exemple, \emph{soit} fournir un raisonnement général.
  \item Concluez en rappelant l’importance de la preuve en mathématiques.
\end{itemize}
\end{exercise}

\begin{exercise}
On définit une fonction $T$ ainsi :

\[
T(n) =
\begin{cases}
\dfrac{n}{2}, & \text{si $n$ est pair}, \\[6pt]
3n + 1, & \text{si $n$ est impair}.
\end{cases}
\]

La \textbf{conjecture de Collatz} dit ceci :
Pour tout entier $n > 1$, si on applique plusieurs fois $T$,
la suite obtenue arrive toujours à $1$, puis continue dans le cycle :

\[
1 \;\;\to\;\; 4 \;\;\to\;\; 2 \;\;\to\;\; 1 \;\;\to\;\; \cdots
\]

\end{exercise}























